% !PRINTED 426
\weberpage{426}
Demnach lautet jetzt die Formel

\weq{10}{\Sigma k\, h(m) = 2\Sigma\partial,}

$n$ kein Quadrat,
\[
= 2\Sigma\partial + \sqrt{n} + \tfrac{1}{6},
\]
$n$ ein Quadrat.

Die Summe auf der linken Seite von (10) zerlegen wir in ihre einzelnen Bestandteile, indem wir jede Lösung der Gleichung (8) besonders nehmen. Dann bekommen wir eine Summe von Ausdrücken $h(m)$, wo aber dasselbe $m$ so oft vorkommt, als (8) Lösungen hat. Eine der Lösungen ergibt aber
\[
m = \frac{4n-y^2}{x^2},
\]
und diese kommt, wenn $y=0$ ist, \gesperrt{einmal}, wenn $y$ von Null verschieden ist, zweimal (mit $+y$ und $-y$) vor. Es ist aber die Anzahl der primitiven Klassen der Diskriminante $-m$ gleich der Anzahl der nicht primitiven Klassen der Diskriminante $-x^2m$ vom Teiler $x$, und wenn wir also jetzt (zum Unterschied von der vorigen Formel) mit $h(m)$ die \gesperrt{Gesamtzahl} der Klassen von der Diskriminante $-m$ (primitiven und imprimitiven) verstehen, so erhalten wir aus (10):
\weq{11}{\begin{aligned}
&h(4n) + 2h(4n-1) + 2h(4n-4) + 2h(4n-9) + \cdots \\
&= 2\sum_{\partial>\sqrt{n}}\partial, \quad (n \text{ kein Quadrat}), \\
&= 2\sum_{\partial>\sqrt{n}}\partial + \sqrt{n} + \tfrac{1}{6}, \quad (n \text{ ein Quadrat}).
\end{aligned}}

Darin ist links die Summe der $h(4n-y^2)$ so lange fortzusetzen, als $4n-y^2$ positiv bleibt, und rechts durchläuft $\partial$ alle Divisoren von $n$, die größer als $\sqrt{n}$ sind.

\gesperrt{Dabei ist jedoch zu beachten, daß die Formen $(x,0,x)$ nur je mit $\frac{1}{2}$ und $(x,x,x)$ mit $\frac{1}{3}$ in Rechnung zu setzen sind,} weil sie aus den Darstellungen von $4n$ durch die Diskriminante $m=-4,-3$ hervorgehen.

\begin{center}
\S\ 117. \textbf{\gesperrt{Arithmetische Natur der Klassenfunktion $H_m(u)$.}}
\end{center}

Wir kehren jetzt wieder zurück zu der Gleichung (10), § 115, um daraus die Natur der Koeffizienten der Klassengleichung $H_m(u)=0$ abzuleiten. Zunächst aber betrachten wir die beiden speziellen Fälle: $m=4$ und $m=3$.
